%======================================================================
\section{AEIC Decomposition and Quality Signals}
\label{sec:aeic}
%======================================================================

We introduce the AEIC decomposition, a structured reasoning protocol that enables reliable conclusion extraction across diverse communication topologies. The quality signals derived from AEIC components additionally serve as a formal interface for Stackelberg mechanism design (Section~\ref{sec:game}).

\subsection{Structured Reasoning Decomposition}

Given a question $x$, each agent $i \in \{1, \ldots, n\}$ produces a structured response decomposed into four components:

\begin{itemize}
    \item \textbf{Assumption (A)}: The premise or hypothesis the agent adopts.
    \item \textbf{Evidence (E)}: A list of factual claims supporting the analysis: $E_i = \{e_{i,1}, \ldots, e_{i,k_i}\}$.
    \item \textbf{Inference (I)}: The reasoning chain connecting evidence to conclusion.
    \item \textbf{Conclusion (C)}: The final answer with justification.
\end{itemize}

This decomposition is enforced through prompt engineering and serves two purposes: (1)~it structures agent output into parseable components, enabling reliable conclusion extraction and consistent stance identification across all communication topologies---this is the primary empirical value (Section~\ref{sec:exp}); and (2)~it provides separable components from which quality signals can be computed, serving as a formal interface for Stackelberg mechanism design (Section~\ref{sec:game}), though experiments reveal that these signals have limited discriminative power when agents share the same LLM backend (Section~\ref{sec:exp}).

\subsection{Four-Dimensional Quality Signals}

From each agent's AEIC output, we compute four quality signals using embedding-based similarity:

\begin{definition}[Quality Signals]
Let $\mathrm{sim}(\cdot, \cdot)$ denote cosine similarity in a sentence embedding space. For agent~$i$:
\begin{align}
    \mathrm{IC}_i &= \mathrm{sim}(I_i, C_i) \\
    \mathrm{EC}_i &= \tfrac{1}{|E_i|}\textstyle\sum_{e \in E_i} \mathrm{sim}(e, C_i) \\
    \mathrm{EQ}_i &= f_{\mathrm{quality}}(E_i) \\
    \mathrm{EV}_i &= \tfrac{1}{|\mathcal{N}_i|}\textstyle\sum_{j \in \mathcal{N}_i} \mathrm{sim}(E_i, E_j)
\end{align}
where IC measures Inference--Conclusion consistency, EC measures Evidence--Conclusion support, EQ assesses evidence quality, and EV captures cross-agent evidence consensus.
\end{definition}

Here $\mathcal{N}_i$ denotes the neighbors of agent $i$ in the communication topology graph, and $f_{\text{quality}}$ assesses evidence coverage and specificity. By design, IC, EC, and EQ measure \emph{individual reasoning coherence}: how internally consistent, well-supported, and thorough an agent's analysis is. EV measures \emph{inter-agent agreement}: how much an agent's evidence aligns with its neighbors'. This distinction matters for incentive design---EV can be inflated by sycophantic behavior. However, as our experiments demonstrate (Section~\ref{sec:exp}), these coherence signals cannot distinguish correct from incorrect reasoning when agents share the same LLM: a coherent wrong answer scores as highly as a correct one.

\subsection{Agent Type Space}
\label{sec:type}

Following Saig et al.~\cite{saig2024incentivizing}, we model each agent through an observable \emph{type}:

\begin{definition}[Agent Type]
The type of agent $i$ on question $q$ is:
\begin{equation}
    \tau_i = (\text{IC}_i, \text{EC}_i, \text{EQ}_i) \in [0,1]^3
\end{equation}
\end{definition}

EV is deliberately excluded from the type. It measures consensus rather than reasoning quality---a sycophantic agent can achieve high EV by mimicking neighbors without genuine reasoning. The type captures what the mechanism is \emph{designed} to reward: individual reasoning effort. In practice, Section~\ref{sec:exp} shows that these type signals have limited discriminative power when all agents share the same LLM backend.

\subsection{AEIC as Feature Extraction}
\label{sec:aeic_features}

Beyond mechanism design, AEIC-structured outputs provide the basis for topology routing features. For each question--topology pair, we construct a 42-dimensional feature vector from per-agent quality signals, summary statistics, \emph{voting features} (majority ratio, yes/no distribution extracted from the Conclusion~component), and cross-signal interactions. Feature ablation (Section~\ref{sec:exp}) reveals that the voting features---enabled by reliable conclusion extraction from AEIC structure---are the primary driver of routing accuracy, while quality signals contribute marginally or negatively.
